\documentclass[12pt,a4paper]{article}
\usepackage[utf8]{inputenc}
\usepackage[T1]{fontenc}
\usepackage[spanish]{babel}
\usepackage[margin=1.4cm]{geometry}
\usepackage{array,tabularx}
\usepackage{enumitem}
\usepackage{titlesec}
\usepackage{xcolor}
\usepackage{hyperref}

% Restrained, high-contrast palette shared by headings, emphasis and links.
\definecolor{HeadingColor}{HTML}{243B53}
\definecolor{InstructionColor}{HTML}{704214}
\definecolor{ConceptColor}{HTML}{205C50}
\definecolor{LinkColor}{HTML}{0047AB}
\hypersetup{colorlinks=true,allcolors=LinkColor,pdfborder={0 0 0}}
\newcommand{\guidelink}[2]{\href{#1}{\textcolor{LinkColor}{\textbf{\underline{#2}}}}}
\newcommand{\instruction}[1]{\textcolor{InstructionColor}{\textbf{#1}}}
\newcommand{\concept}[1]{\textcolor{ConceptColor}{\textbf{#1}}}

\setlength{\parindent}{0pt}
\setlength{\parskip}{3pt}
\setlength{\tabcolsep}{5pt}
\renewcommand{\arraystretch}{1.12}
\setlist{leftmargin=0.6cm,topsep=3pt,partopsep=0pt,parsep=0pt,itemsep=1pt}
\titleformat{\section}{\normalsize\bfseries\color{HeadingColor}}{}{0pt}{}
\titlespacing*{\section}{0pt}{8pt}{3pt}
\newcolumntype{Y}{>{\raggedright\arraybackslash}X}
\newcommand{\checkbox}{\fbox{\rule{0pt}{0.8ex}\rule{0.8ex}{0pt}}}

\begin{document}

\begin{center}
{\Large\bfseries\color{HeadingColor} C9--C10: Explorar ideas desde la comprensión del usuario}\\[3pt]
Emprendimiento e Innovación --- Grado 10
\end{center}

\section*{Objetivo y punto de partida}
\concept{C9:} Participar en la ideación inicial con apertura a propuestas diversas.\\
\concept{C10:} Explicar por qué comprender al usuario debe orientar y limitar las propuestas.

Ya cuentan con los hallazgos H de C5, la oportunidad priorizada en C6, el desafío D1 de C7 y las conexiones K1/K2 de C8. C9 explora posibilidades a partir de estos resultados; C10 examina sus supuestos. \instruction{Las ideas son exploratorias: no elegirán una solución, ni la desarrollarán, prototiparán o implementarán.}

\section*{Documento acumulativo e integridad de la evidencia}
Esta guía es de consulta. Añadan C9 y C10 al \textbf{mismo documento del proyecto en LaTeX}; conserven las secciones previas y actualicen la fecha. Pueden continuar en \guidelink{https://www.overleaf.com/}{Overleaf}. Usen referencias internas, sin copiar D1, perfiles, mapas, hallazgos ni matrices. Entreguen el PDF completo actualizado.

Mantengan la distinción entre evidencia, conclusión, supuesto e idea. La IA solo puede apoyar código, formato y organización de contenido propio; \instruction{no debe fabricar evidencia, decisiones, ideas ni contribuciones de integrantes.} Los ejemplos son ficticios. Mantengan identificada cualquier evidencia simulada autorizada.

\section*{C9 --- Preparen y abran la exploración}
Lean D1, los hallazgos H pertinentes y K1/K2. Identifiquen oralmente la necesidad respaldada y un supuesto pendiente. No entreguen otro resumen. Si detectan una inconsistencia, corríjanla en la sección de origen y dejen allí una nota del cambio.

Durante \textbf{cinco minutos}, cada integrante escriba \textbf{tres ideas diferentes}, sin consultar ni evaluar las de otros. Cada idea debe expresar qué podría hacerse y cómo intentaría atender D1. Exploren distintas formas de organizar, comunicar, actuar o facilitar la experiencia; cambiar el nombre de una propuesta no crea una alternativa distinta.

Asignen códigos I1, I2, etc. y registren directamente las ideas en la tabla única de C9. No creen primero un listado para volver a transcribirlo. No necesitan dibujos, costos ni detalles técnicos.

\section*{C9 --- Escuchen, amplíen y registren}
\begin{enumerate}
\item Por turnos, cada integrante presenta sus ideas. Los demás escuchan sin interrumpir ni descartarlas. Cada integrante hace al menos una pregunta para comprender una propuesta ajena.
\item Cada integrante aporta una ampliación, adaptación o combinación de una idea de otra persona. Utilicen K1/K2 para abrir perspectivas; no es obligatorio que cada idea aplique ambos conceptos.
\item Añadan \textbf{dos ideas colaborativas} nuevas o modificadas a la misma tabla, con códigos propios y referencias a sus ideas de origen. Conserven todas las propuestas, incluso las que no combinaron. No voten ni ordenen ganadoras.
\end{enumerate}

{\small
\begin{tabularx}{\linewidth}{|p{0.08\linewidth}|Y|Y|Y|}
\hline
\textbf{Código} & \textbf{Idea: descripción breve} & \textbf{Relación con D1} & \textbf{Autoría y construcción} \\
\hline
I1, etc. & Posibilidad por explorar. & Necesidad y hallazgo H; aporte de K1/K2, si corresponde. & Integrante; en ideas colaborativas, códigos de origen, participantes y cambio realizado. \\
\hline
\end{tabularx}}

Para evidenciar la apertura, añadan en la última columna de la \textbf{idea ajena correspondiente} una nota por integrante: quién escuchó, qué preguntó y qué aportó. No elaboren otro registro de asistencia o participación. Si varias personas comentan la misma idea, distingan sus aportes.

\concept{Evidencia de C9:} un registro de todas las ideas, tres por integrante y dos colaborativas, con trazabilidad de autoría, escucha y ampliación. \textbf{Una idea es una posibilidad, no evidencia de que funcionará.}

\newpage
\section*{C10 --- Examinen cómo el usuario condiciona las ideas}
Elijan \textbf{dos ideas distintas para analizar}, al menos una colaborativa. Esta muestra no constituye una selección de finalistas. Referencien sus códigos y los registros del usuario de C5; no reconstruyan perfiles ni realicen nuevas entrevistas para esta entrega.

Para cada idea, completen una fila. Cada columna cumple una función distinta:
\begin{itemize}
\item \textbf{Evidencia:} qué conocen del usuario, con hallazgo H y fuente existente.
\item \textbf{Supuesto y riesgo:} qué están dando por cierto y qué podría ocurrir si solo se guían por sus preferencias.
\item \textbf{Condición y consecuencia:} qué debe respetar la idea según la evidencia; ajuste necesario o razón concreta para mantenerla.
\item \textbf{Pendiente:} qué necesitan conocer del usuario y cómo podrían averiguarlo en una etapa posterior, sin pedirle que confirme su idea.
\end{itemize}

{\small
\begin{tabularx}{\linewidth}{|p{0.08\linewidth}|Y|Y|Y|Y|}
\hline
\textbf{Idea} & \textbf{Evidencia del usuario} & \textbf{Supuesto y riesgo} & \textbf{Condición y consecuencia} & \textbf{Pendiente} \\
\hline
Código I & Hallazgo H y fuente. & Supuesto no verificado; efecto posible de ignorar al usuario. & Límite respaldado; ajuste o razón para mantener la idea. & Pregunta neutral y forma futura de aclararla. \\
\hline
\end{tabularx}}

Si deben ajustar una idea, editen su fila de C9 y anoten aquí el cambio y el hallazgo que lo motiva; no copien la propuesta completa. Si no hay respaldo para una condición, reconózcanlo: no inventen evidencia. No puntúen ni comparen ganadoras.

\textbf{Ejemplo ficticio:} si un registro indica que algunos usuarios no disponen de teléfono durante el descanso, una idea que lo exige puede excluirlos. Esa evidencia condiciona la propuesta; no demuestra que otra alternativa funcionará. Antes de desarrollarla habría que comprender cómo acceden actualmente al servicio.

\section*{C10 --- Expliciten el aprendizaje}
Cierren el mismo análisis con \textbf{dos frases de síntesis}, sin repetir las filas: expliquen por qué la evidencia del usuario debe orientar las ideas desde el inicio y por qué conocer una necesidad no demuestra la aceptación ni el funcionamiento de una propuesta. Apóyense en el contraste realizado, no en una afirmación general como ``hay que escuchar al usuario''.

\concept{Evidencia de C10:} dos contrastes que hacen visibles riesgos, condiciones y preguntas pendientes, más una explicación de la importancia de comprender al usuario. C10 aplica la evidencia de C2--C5; no repite su recopilación ni vuelve a priorizar oportunidades.

\section*{Resultado y continuidad}
El documento reúne ahora una secuencia completa: contexto organizado (C5), oportunidad priorizada (C6), desafío definido (C7), conexiones disciplinares (C8), ideación abierta (C9) y límites derivados del usuario (C10). Las propuestas permanecen abiertas. Los pendientes orientarán una investigación posterior; todavía no existe una solución elegida o validada.

\section*{Verificación final}
{\small\begin{itemize}[label=\checkbox]
\item Usamos D1, H y K1/K2 sin recrear los productos anteriores.
\item Cada integrante generó tres ideas y documentó una pregunta y un aporte a una idea ajena.
\item Una sola tabla conserva todas las ideas y dos propuestas colaborativas con sus orígenes.
\item C10 distingue evidencia, supuestos, riesgos, condiciones y pendientes en dos ideas.
\item Explicamos cómo comprender al usuario orienta las propuestas, sin afirmar que ya funcionan.
\item No elegimos soluciones ni fabricamos aportes; compilamos y revisamos el PDF acumulativo.
\end{itemize}}

\end{document}
